\documentclass[10pt,a4paper]{amsart}
\usepackage[margin=19mm]{geometry}
\usepackage{amsmath,amssymb,mathtools}
\usepackage[scr=rsfs,cal=euler]{mathalfa}
\usepackage{xfrac}
\usepackage[unicode,hidelinks,bookmarks=false]{hyperref}
\newcommand{\ie}{i.e.\ }
\newcommand{\cf}{cf.\ }
\newcommand{\resp}{respectively\ }
\newcommand{\Subsection}{\subsection}
\providecommand{\nobreakdash}{\nobreak\hbox{-}}
\setlength{\emergencystretch}{2em}
\multlinegap0pt
\usepackage{amssymb}
\usepackage{mathtools}
\usepackage{xcolor}
\usepackage{tikz}
\newtheorem{lemma}{Lemma}
\newtheorem{prop}{Proposition}
\newcommand{\R}{\mathbb{R}}
\newcommand{\cT}{\mathcal{T}}
\newcommand{\cU}{\mathcal{U}}
\newcommand{\cW}{\mathcal{W}}
\DeclareMathOperator{\diam}{diam}
\newcommand{\from}{\colon}
\DeclareMathOperator{\gen}{gen}
\DeclareMathOperator{\mesh}{mesh}
\title{Area of H\"older curves and coarea formula on the Heisenberg group}
\author{Gioacchino Antonelli, Robert Young}
\date{Source: arXiv:2605.15987v1 (CC BY 4.0)}
\begin{document}
\maketitle

\noindent\emph{Source and licence.} Area of H\"older curves and coarea formula on the Heisenberg group. Version arXiv:2605.15987v1 (CC BY 4.0), distributed by arXiv under the Creative Commons Attribution 4.0 International licence, http://creativecommons.org/licenses/by/4.0/, as recorded in the arXiv licence metadata for this exact version and shown on the abstract page https://arxiv.org/abs/2605.15987v1. Full licence notice: \texttt{licenses/arxiv-2605.15987-CC-BY-4.0.txt}.\par\smallskip

\noindent\textit{Editorial scope.} Counted unit: the article's prop prop:rn-curve-area (source0.tex lines 1997--2022). The article's complete original proof of it is delivered with it (source0.tex lines 2039--2166, one \texttt{proof} environment of 128 lines, which the article places away from the statement). No mathematics is written, repaired or re-derived: the statement and the proof are the article's own lines, in the article's own notation, and no retained line is substituted or re-flowed.  Retained uncounted as the source-local context the proof needs: lines 538--553; lines 961--970; lines 1400--1400; lines 1417--1427; lines 1419--1421; lines 1423--1425; lines 1471--1486; lines 1479--1481; lines 1539--1541; lines 1554--1574; lines 1560--1562; lines 1579--1581; lines 1584--1591; lines 1607--1612; lines 1629--1631; lines 1644--1646; lines 1693--1693; lines 1786--1786; lines 1813--1815.  Nothing was omitted from any retained span, so the package carries no omission record.  No retained line contains a citation, so the wrapper carries no bibliography and no bibliography entry.  No retained line contains a figure environment, an image-inclusion command or drawing code, so no image is delivered and the package declares no asset dependency.  Editorial formatting only, in addition: an AMS \texttt{amsart} preamble that restates the article's own declarations and macros used by the retained lines, namely the environments \texttt{lemma}, \texttt{prop} on separate counters, together with the standard packages those lines need.  The retained environments print in this document as Lemma 1, Proposition 1 in order of appearance, whereas the article prints the same items with the numbering of its own full document.  The complete native source of the article as distributed in the arXiv e-print for this exact version is bundled unchanged as \texttt{sources/200696/source0.tex}.
\begin{lemma}\label{lem:symplectic-area-props-v2}   Let $I\subset \R$ be a closed interval. The following properties hold:
  \begin{enumerate}
  \item Let $I,J\subset \R$ be compact intervals and let $\gamma\from I\to\mathbb R^{2n}$ be a continuous curve. If $S(\gamma)$ exists and if $\alpha\from J\to I$ is a monotone nondecreasing homeomorphism, then $S(\gamma\circ \alpha) = S(\gamma)$.
  \item If $\gamma_1$ and $\gamma_2$ are continuous curves, $S(\gamma_1)$ and $S(\gamma_2)$ exist, and $\gamma_1\diamond \gamma_2$ is their concatenation, then $S(\gamma_1\diamond \gamma_2) = S(\gamma_1)+S(\gamma_2)$.
  \item Let $\gamma\from I\to\mathbb R^{2n}$ be a closed curve with finite length. Then 
  \[
  |S(\gamma)|\le \ell(\gamma)^2,
  \]
  where $\ell(\gamma)$ is the length of $\gamma$.
  \item Let $\gamma,\gamma_i\from I\to\mathbb R^{2n}$ be curves with finite length, and assume that $\gamma_i\to\gamma$ uniformly on $I$, and that $\sup_{i\in\mathbb N}\ell(\gamma_i)<\infty$. Then the symplectic area of $\gamma$ exists and
  \[
  \lim_{i\to\infty}S(\gamma_i)=S(\gamma).
  \]
\item If $\alpha\from I\to \R^{2n}$ is a closed curve, $v\in \R^{2n}$, and $\beta(t) = \alpha(t) + v$, then $S(\alpha) = S(\beta)$.
  \end{enumerate}
\end{lemma}

\begin{prop}\label{prop:vertical-curve-area}
  With notation as above, if $\sigma(\Lambda)<\infty$, then
  $S(\hat{\gamma})$ exists and 
  $S(\hat{\gamma}) = \lim_{i\to \infty} S(\hat{g}_i)$. Furthermore, 
  \begin{equation}\label{eq:sgamma-bound}
    |S(\hat{\gamma})| \lesssim_{\mu,C} \max\left\{\sigma(\Lambda), \sigma(\Lambda) \log \frac{(\diam K)^2}{\sigma(\Lambda)}\right\}.
  \end{equation}
  
  If, in addition, $g_i(0)=g_i(1)$ for all $i$, then $|S(\hat{\gamma})| \lesssim_{\mu,C} \sigma(\Lambda)$.
\end{prop}

\section{Area formula based on approximations}\label{sec:area-formula}

\begin{lemma}\label{lem:nbhd-diam}
  For every $v\in \cT^i$,
  \begin{equation}\label{eq:delta-lambda-bound}
    \delta_\Lambda(v) \lesssim 2^{-i}\diam K.
  \end{equation}
  For every $(t,u) \in [0,1]\times [0,1]$,
  \begin{equation}\label{eq:G-gamma-dist}
    |G(t,u) - \gamma(t)| \lesssim u \diam(K),
  \end{equation}
  so $G$ is continuous.
\end{lemma}

  \begin{equation}
    \delta_\Lambda(v) \lesssim 2^{-i}\diam K.
  \end{equation}

  \begin{equation}
    |G(t,u) - \gamma(t)| \lesssim u \diam(K),
  \end{equation}

\begin{lemma}\label{lem:sgi-converge}
  For any $v\in \cT$, 
  \begin{equation}\label{eq:ell-theta-v}
    \ell(\theta_v) \lesssim \delta_\Lambda(v) \lesssim 2^{-\gen(v)}\diam(K).
  \end{equation}
  Therefore,
  $$\sum_{v\in \cT} |S(\theta_v)| \lesssim \sum_{v\in \cT} \delta_\Lambda(v)^2 = \sigma(\Lambda)<\infty.$$
  It follows that $\sum_{v\in \cT} S(\theta_v)$ converges absolutely. Let
  \begin{equation}\label{eq:def-psi}
    \psi = (G\circ \overline{(1,0),(1,1)}) \diamond (G\circ \overline{(0,1),(0,0)}).
  \end{equation}
  Since $G$ is constant on the top edge of $[0,1]^2$, this is a curve from $G(1,0) = \gamma(1)$ to $G(0,0) = \gamma(0)$.
  Let $\Phi = \sum_{v\in \cT} S(\theta_v)$. 
  Then $\ell(\psi)<\infty$ and
  $$\lim_{i\to \infty} S(g_i) = \Phi - S(\psi).$$
\end{lemma}

  \begin{equation}
    \psi = (G\circ \overline{(1,0),(1,1)}) \diamond (G\circ \overline{(0,1),(0,0)}).
  \end{equation}

\begin{lemma}\label{lem:boundary-arches}
  The set $R(\cU_P)$ is a simply-connected noncompact planar manifold, homeomorphic to a closed disc with $k+1$ boundary points removed. For each $i=1,\dots,k$, $R(\cU_P)$ has a boundary component connecting $(t_{i-1},0)$ to $(t_{i},0)$, and this component can be parameterized by a map $u=(\alpha_1,\alpha_2)\from [0,1]\to \partial R(\cU_P)$ such that $\alpha_1$ is nondecreasing and $\alpha_2(x)>0$ for all $x\in (0,1)$.
\end{lemma}

\begin{lemma}\label{lem:arch-lengths}
  For each $j$, let
  \begin{equation}\label{eq:def-cwj}
    \cW_j = \{w\in \cT : J_w\cap \{t_{j-1},t_j\} \ne \emptyset \text{ and } \diam I_w \le 2 \mu^2 \diam \zeta([t_{j-1},t_j])\}.
  \end{equation}
  Then 
  \begin{equation}\label{eq:boundary-graph}
    A_j\setminus \{(t_{j-1},0), (t_{j},0)\} \subset \bigcup_{w\in \cW_j} \partial R_w
  \end{equation}
  and 
  \begin{equation}\label{eq:alpha-length}
    \ell(a_j) \lesssim \sum_{w\in \cW_j} \delta_\Lambda(w) \lesssim  d(\zeta(t_{j-1}),\zeta(t_j)).
  \end{equation}
  Furthermore, the $\cW_j$'s satisfy the following.
  \begin{enumerate}
  \item\label{it:W-multiplicity}
    There is a $c = c(\mu,n) >0$ such that for any $v\in \cT$, $|\{j : v\in \cW_j\}|\le c$.
  \item\label{it:arch-count}
    For all $i\ge 0$ and all $j$, $|\cW_j \cap \cT^i| \le 4$.
  \end{enumerate}
\end{lemma}

  \begin{equation}
    A_j\setminus \{(t_{j-1},0), (t_{j},0)\} \subset \bigcup_{w\in \cW_j} \partial R_w
  \end{equation}

\begin{equation}\label{eqn:RelationPhiPAlphaP}
\phi_P = G|_{\partial R(\cU_P)} = \alpha_P \diamond \psi,
\end{equation}

\begin{lemma}\label{lem:lim-dyadic}
  For any partition $P$, we have $|S(\phi_P)| \lesssim \sigma(\Lambda)$ and
  \begin{equation}\label{eq:lim-dyadic}
    \lim_{\mesh(P)\to 0} S(\phi_P) = \Phi.
  \end{equation}
  Therefore, by \eqref{eqn:RelationPhiPAlphaP} and Lemma \ref{lem:symplectic-area-props-v2}.(2),
  $$\lim_{\mesh(P)\to 0} S(\alpha_P) = \lim_{\mesh(P)\to 0} (S(\phi_P) - S(\psi)) = \Phi - S(\psi).$$
\end{lemma}

\begin{lemma}\label{lem:arch-area}
  For any $\epsilon>0$, there is a $\rho$ such that if $\mesh P<\rho$, then
  \begin{equation}\label{eq:arch-area}
    \sum_j \ell(a_j)^2 \le \epsilon + \epsilon \sum_j \Delta z_j.
  \end{equation}
\end{lemma}

  \begin{equation}\label{eq:vca-pf-s-gamma}
    S(\gamma_P) = S(\alpha_P) - \sum_j S(\hat{a}_j).
  \end{equation}

  \begin{equation}\label{eq:s-gamma-s-psi}
    S(\hat{\gamma}) = \Phi - S(\psi) + S(\overline{\gamma(1),\gamma(0)}) = \Phi - S(\hat{\psi}).
  \end{equation}

\subsection{Proof of Lemma~\ref{lem:arch-lengths}}\label{sec:proof-arch-lengths}

\subsection{Proof of Lemma~\ref{lem:lim-dyadic}}\label{sec:proof-dyadic}

  \begin{equation}\label{eq:area-phi-P-sum}
    \sum_{v\in \cU_P} \ell(\theta_v) \lesssim \sum_{v\in \cU_P} 2^{-\gen(v)} \diam(K) \le \sum_{i=0}^\infty 2^{-i} \diam(K) |\cU_P\cap \cT^i|.
  \end{equation}

\begin{prop}\label{prop:rn-curve-area}
  Let $L>0$. Let $\gamma\from [0,1] \to \R^{2n}$ be a curve such that
  \begin{equation}\label{eq:rn-gamma-holder}
    |\gamma(s) - \gamma(t)| \le L \sqrt{|s-t|},
  \end{equation}
  for all $s,t\in [0,1]$. For $i\ge 0$ and $j=0,\dots, 2^i$, let $\Lambda_{i,j} \in \R^{2n}$ be such that
  \begin{equation}\label{eqn:Lambdaij}
  |\Lambda_{i,j} - \gamma(j2^{-i})| \leq L 2^{-\frac{i}{2}}.
  \end{equation}
  For each $i$, let $h_i$ be the curve
  \begin{equation}\label{eqn:girn}
  h_i =\overline{\Lambda_{i,0},\dots, \Lambda_{i,2^i}},
  \end{equation}
  that connects $\Lambda_{i,0},\dots, \Lambda_{i,2^i}$ by straight line segments. Let
  $$\delta_{i,j} = \diam(\{\Lambda_{i,j}, \Lambda_{i,j+1}, \Lambda_{i+1,2j}, \Lambda_{i+1,2j+1}, \Lambda_{i+1,2j+2}\}).$$
  If
  \begin{equation}\label{eqn:ConditionSigmaRn}
  \sigma := \sum_{i=0}^\infty \sum_{j=0}^{2^i-1} \delta_{i,j}^2 < \infty,
  \end{equation}
  then $S(\hat{\gamma})$ exists, $S(\hat{\gamma}) = \lim_iS(\hat{h}_i)$, and 
  \begin{equation}\label{eq:sgi-bound-rn}
    |S(\hat{\gamma})| \lesssim \sigma \max\left\{1, \log \frac{L^2}{\sigma}\right\}.
  \end{equation}
  Here $\hat\gamma,\hat h_i$ are the closed curves obtained by joining the endpoints of $\gamma,h_i$ with straight line segments. If $h_i(0)=h_i(1)$ for all $i$, or if $h_i(0) = \gamma(0)$ and $h_i(1) = \gamma(1)$ for all $i$, then 
  $|S(\hat{\gamma})| \lesssim \sigma$.
\end{prop}

\begin{proof}[Proof of Proposition~\ref{prop:rn-curve-area}]
  For all $i$ and $j$, let $\delta_\Lambda(v_{i,j}) = \delta_{i,j}$. The argument that proves Lemma~\ref{lem:nbhd-diam} implies
  \begin{equation}\label{eqn:EstimateDeltaLambdarn}
    \delta_\Lambda(v) \lesssim L 2^{-\frac{i}{2}}
  \end{equation}
  for all $v\in \cT^i$ and   
  \begin{equation}\label{eqn:Gtu}
    |G(t,u)-\gamma(t)|\lesssim L \sqrt{u}
  \end{equation}
  for all $t$ and $u$.  These bounds differ slightly from \eqref{eq:delta-lambda-bound} and \eqref{eq:G-gamma-dist} because those inequalities deal with patchworks whose level--$i$ pieces have diameter on the order of $2^{-i}$, whereas in this situation, $\gamma(J_v)$ can have diameter of order $2^{-\frac{\gen(v)}{2}}$. These estimates will still suffice for the argument.

  Let $\Phi = \sum_{v\in\cT} S(\theta_v)$. As noted above, this sum converges absolutely, so the proof of Lemma~\ref{lem:sgi-converge} shows that
  \[
    \lim_{i\to\infty} S(h_i)=\Phi-S(\psi),
  \]
  where
  $$\psi = (G\circ \overline{(1,0),(1,1)}) \diamond \overline{h_0(1), h_0(0)}\diamond (G\circ \overline{(0,1),(0,0)})$$
  traces out three sides of $G([0,1]^2)$.
  (This is slightly different from \eqref{eq:def-psi} because $h_0$ is not necessarily constant.)
  It remains to show that $S(\gamma) = \lim_{i\to \infty} S(h_i)$.

  For any partition $P=\{t_0,\dots,t_k\}$ of $[0,1]$, we define $\gamma_P$, $\cU_P$, and $R(\cU_P)$ as in Section~\ref{sec:area-formula}.
  Lemma~\ref{lem:boundary-arches} holds as before, so $\partial R(\cU_P)$ contains arches $A_j$ connecting $(t_{j-1},0)$ and $(t_j,0)$. We let $u_j$ parameterize $A_j$, let $a_j = G\circ u_j$, and let $\alpha_P = a_1\diamond \dots \diamond a_k$. Then, as in \eqref{eq:vca-pf-s-gamma}, 
  \begin{equation}\label{eq:s-gamma-s-alpha-rn}
    S(\gamma_P) = S(\alpha_P) - \sum_j S(\hat{a}_j),
  \end{equation}
  so it suffices to estimate $S(\alpha_P)$ and the $\ell(a_j)$'s.

  The arguments in Section~\ref{sec:proof-arch-lengths} show that if
  \begin{equation}\label{eq:def-cwjrn}
    \cW_j := \{w\in \cT : J_w\cap \{t_{j-1},t_j\} \ne \emptyset\, \text{ and}\, \diam J_w \leq 2|t_j-t_{j-1}|\},
  \end{equation}
  then \eqref{eq:boundary-graph} holds and 
  \begin{equation}\label{eq:alpha-length-rn}
    \ell(a_j) \lesssim \sum_{w\in \cW_j} \delta_\Lambda(w) \lesssim \sqrt{t_{j}- t_{j-1}}
  \end{equation}
  for all $j$, as in Lemma~\ref{lem:arch-lengths}. Furthermore, parts (\ref{it:W-multiplicity}) and (\ref{it:arch-count}) of Lemma~\ref{lem:arch-lengths} hold verbatim.

  Likewise, the proof of Lemma~\ref{lem:lim-dyadic} (see Section~\ref{sec:proof-dyadic}) still holds, using the inequality $\ell(\theta_v)\lesssim 2^{-\frac{\gen(v)}{2}}$ rather than $\ell(\theta_v)\lesssim 2^{-\gen(v)}$ in \eqref{eq:area-phi-P-sum}. That is,
  \begin{equation}\label{eq:lim-dyadic-conclusion}
    \lim_{\mesh(P)\to 0} S(\alpha_P) = \lim_{\mesh(P)\to 0} (S(\phi_P) - S(\psi)) = \Phi - S(\psi).
  \end{equation}
  
  Furthermore, we can use the argument of Lemma~\ref{lem:arch-area} to prove that if $\mesh(P)$ is sufficiently small, then
  \begin{equation}\label{eqn:ajepsilon}
    \sum_j \ell(a_j)^2\leq \epsilon.
  \end{equation}

  We proceed as follows. Let $\delta > 0$ and $m\in \mathbb{N}$ be numbers to be chosen later. Let $i_0>0$ be such that
  $$
  \sum_{v\in \cT^{\ge i_0}}\delta_\Lambda(v)^2 \le \delta.
  $$
  Since $\sigma<\infty$, such an $i_0$ exists.
  Let $\rho:=2^{-i_0-2}$.

  Suppose that $\mesh P < \rho$.
  Let $\cW_j\subset \cT$ be as in \eqref{eq:def-cwjrn} and let $n_j = \min_{w\in \cW_j} \gen(w)$. For any $w\in \cW_j$,
  \begin{equation}\label{eqn:SqrtEstimate}
  2^{-n_j}=\diam J_w \leq 2|t_j-t_{j-1}|\leq 2\mesh P < 2^{-i_0},
  \end{equation}
  so $n_j > i_0$. Thus, by \eqref{eq:alpha-length-rn},
  \[
  \begin{aligned}
    \ell(a_j)^2 \lesssim \left(\sum_{w\in \cW_j} \delta_\Lambda(w)\right)^2  
    \lesssim \left(\sum_{w\in \cW_j(<n_j + m)}\delta_\Lambda(w)\right)^2 + \left(\sum_{w\in \cW_j(\ge n_j + m)}\delta_\Lambda(w)\right)^2,
\end{aligned}
  \]
  where $\cW_j(< k) := \cW_j\cap \cT^{< k}$ and $\cW_j(\ge k) := \cW_j\cap \cT^{\ge k}$.

  We bound these terms separately. On one hand, as in Lemma~\ref{lem:arch-lengths}.(\ref{it:arch-count}), we have $|\cW_j\cap\cT^k|\le 4$ for all $k$, so $|\cW_j(<n_j + m)| \le 4m$. Therefore,
  $$
  \left(\sum_{w\in \cW_j(<n_j + m)}\delta_\Lambda(w)\right)^2 \le 4m \sum_{w\in \cW_j(<n_j + m)}\delta_\Lambda(w)^2,
  $$
  by Cauchy--Schwarz. Likewise,
  \begin{multline*}
    \sum_{w\in \cW_j(\ge n_j + m)}\delta_\Lambda(w)
    \stackrel{\eqref{eqn:EstimateDeltaLambdarn}}{\lesssim} L \sum_{k=n_j+m}^\infty |\cW_j\cap\cT^k| 2^{-\frac{k}{2}} \\
    \lesssim L \cdot 2^{-\frac{m}{2}} 2^{-\frac{n_j}{2}} \stackrel{\eqref{eqn:SqrtEstimate}}{\lesssim} L2^{-\frac{m}{2}}\sqrt{|t_j-t_{j-1}|}.
  \end{multline*}
  That is,
  $$\ell(a_j)^2\lesssim 4m \sum_{w\in \cW_j(<n_j + m)}\delta_\Lambda(w)^2 + L^22^{-m} |t_j-t_{j-1}|,$$
  and
  $$\sum_j \ell(a_j)^2\lesssim 4m \sum_j \sum_{w\in \cW_j(<n_j + m)}\delta_\Lambda(w)^2 + L^22^{-m}.$$

  Note that $\bigcup_j \cW_j \subset \cT^{>i_0}$ and that each $v\in \cT$ appears in at most $c$ of the $\cW_j$'s (as in Lemma~\ref{lem:arch-lengths}.(\ref{it:W-multiplicity})). Therefore, 
  $$
  \sum_j \sum_{w\in \cW_j(<n_j + m)}\delta_\Lambda(w)^2 \le \sum_{w\in \bigcup_j \cW_j}\delta_\Lambda(w)^2 \le c \sum_{w\in \cT^{>i_0}}\delta_\Lambda(w)^2 \le c\delta,
  $$
  and
  $$\sum_j \ell(a_j)^2\lesssim 4cm\delta + L^22^{-m}.$$
  If we choose $m$ large enough and $\delta$ small enough, this implies \eqref{eqn:ajepsilon}. Thus, by Lemma~\ref{lem:symplectic-area-props-v2}.(3), 
  \begin{equation}\label{eq:s-hat-aj-rn}
    \lim_{\mesh P \to 0}\sum_j S(\hat{a}_j) = 0.
  \end{equation}

  Finally,
  \begin{equation}\label{eq:sgamma-rn}
    S(\gamma) = \lim_{\mesh P \to 0} S(\gamma_P) \stackrel{\eqref{eq:s-gamma-s-alpha-rn}}{=} \lim_{\mesh P \to 0} S(\alpha_P) - \sum_j S(\hat{a}_j)  \stackrel{\eqref{eq:lim-dyadic-conclusion} \wedge \eqref{eq:s-hat-aj-rn}}{=} \Phi - S(\psi).
  \end{equation}

  As in \eqref{eq:s-gamma-s-psi}, we have $S(\hat\gamma)=\Phi-S(\hat\psi)$ and $|\Phi|\lesssim\sigma$. It remains to show that 
  \begin{equation}\label{eq:rn-target}
    S(\hat{\psi}) \lesssim \ell(\psi)^2 \lesssim \sigma\max\left\{1,\log\frac{L^2}{\sigma}\right\}.
  \end{equation}
  This follows from essentially the same argument as in the proof of Proposition~\ref{prop:vertical-curve-area}, but there are a few changes; we briefly sketch the argument here.

  For $t=0,1$, let $\psi_t = G\circ \overline{(t,1),(t,0)}$ and let $\tilde{\psi}_1$ be the reverse of $\psi_1$, so that $\psi = \tilde{\psi}_1 \diamond \overline{h_0(1), h_0(0)}\diamond \psi_0$. 
  Then $\psi_t$ consists of segments from $h_i(t)$ to $h_{i+1}(t)$. By \eqref{eqn:Gtu}, we have $|h_{i+1}(t)-h_i(t)|\lesssim L2^{-\frac{i}{2}}$, and by \eqref{eq:rn-gamma-holder} and \eqref{eqn:Lambdaij}, we have $|h_0(0)-h_0(1)| \lesssim L$, so
  \begin{equation}\label{eq:ell-psi-rn}
    \ell(\psi) \lesssim |h_0(0)-h_0(1)| + \sum_{i=0}^\infty L2^{-\frac{i}{2}} \lesssim L.
  \end{equation}

  If $\sigma > \frac{1}{16} L^2$, then $\ell(\psi)^2 \lesssim L^2 \lesssim \sigma$, implying \eqref{eq:rn-target}. Otherwise, we let
  $$i_0=\left\lfloor \log_2 \frac{L^2}{\sigma}\right\rfloor;$$
  note that $i_0 \ge 4$.  Then $|h_0(0)-h_0(1)| \lesssim \delta_{\Lambda}(v_{0,0})$, $|h_{i+1}(0)-h_i(0)|\le \delta_{\Lambda}(v_{i,0})$, and $|h_{i+1}(1)-h_i(1)|\le \delta_{\Lambda}(v_{i,2^i-1})$, so we can sharpen \eqref{eq:ell-psi-rn} to
  \begin{multline*}
    \ell(\psi) \lesssim \sum_{i=0}^{i_0-1} \left(\delta_\Lambda(v_{i,0}) + \delta_\Lambda(v_{i,2^i})\right) + \sum_{i=i_0}^{\infty} L2^{-\frac{i}{2}}\\
    \lesssim \sum_{i=0}^{i_0-1} \left(\delta_\Lambda(v_{i,0}) + \delta_\Lambda(v_{i,2^i})\right) + L2^{-\frac{i_0}{2}}.
  \end{multline*}
  By Cauchy--Schwarz,
  $$\ell(\psi)^2 \lesssim i_0 \bigg(\sum_{i=0}^{i_0-1} \left(\delta_\Lambda(v_{i,0}) + \delta_\Lambda(v_{i,2^i})\right)^2 + L^22^{-i_0}\bigg) \lesssim
  i_0 (\sigma + L^22^{-i_0}) \lesssim i_0 \sigma,$$
  as desired.

  Finally, if $h_i(0)=h_i(1)$ for all $i$, then $\psi_0 = \psi_1$, and $S(\psi) = S(\tilde{\psi}_1 \diamond \psi_0) = 0$. If $h_i(0) = \gamma(0)$ and $h_i(1) = \gamma(1)$ for all $i$, then $\psi = \overline{\gamma(1),\gamma(0)}$ and $\ell(\psi)^2 \le \delta_\Lambda(v_{0,0})^2 \le \sigma$. In either case, $S(\hat{\psi})\lesssim \sigma$, so \eqref{eq:sgamma-rn} implies that
  $$|S(\hat\gamma)| \le |\Phi| + |S(\hat\psi)| \lesssim \sigma,$$
  as desired.
\end{proof}

\end{document}
